\documentclass[11pt]{article}

\usepackage[margin=1in]{geometry}
\usepackage{amsmath, amssymb, amsthm}
\usepackage{graphicx}
\usepackage{hyperref}
\usepackage{enumitem}
\usepackage{xcolor}

\title{
    \normalsize IDs: 213851579, 213310485 \\[1em] 
    \huge Video Processing -- HW3 \\
    \large Particle Filter Tracking -- Part A
}

\author{}

\date{Course 0512.4263 \quad 2025/2026}

\begin{document}

\maketitle

% ========================================================================
\section*{Part A -- Open Questions}
% ========================================================================

% ------------------------------------------------------------------------
\subsection*{Q1.}
% ------------------------------------------------------------------------
\noindent\textit{We saw that $\text{posterior} \propto \text{prior} \times \text{likelihood}$. Explain the formula
and how it relates to the Kalman/Particle filter. Start from explaining what is the posterior,
prior and likelihood.}

\medskip
\noindent\textbf{Answer.}

By Bayes' rule, $\text{posterior} = (\text{likelihood} \times \text{prior}) / \text{evidence}$. The evidence $p(z)$ does
not depend on the state, so it is only a normalizing constant and we drop it, giving
$\text{posterior} \propto \text{prior} \times \text{likelihood}$.

\begin{itemize}[nosep]
    \item \textbf{Prior} $p(x_t \mid z_{1:t-1})$: the belief about the state before seeing the
    new frame. We get it by pushing the previous estimate through the motion model (the
    predict step).
    \item \textbf{Likelihood} $p(z_t \mid x_t)$: how well a candidate state fits the new
    measurement. In our tracker this is the histogram similarity between the candidate's
    patch and the reference template $q$.
    \item \textbf{Posterior} $p(x_t \mid z_{1:t})$: the belief after using the new frame. It is
    large for states that the prior expected and the likelihood supports.
\end{itemize}

Both filters compute this. Kalman assumes every distribution is Gaussian and updates the mean
and covariance analytically. The particle filter drops that assumption: the prior is the set of
predicted particles, the likelihood becomes each particle's weight, and the posterior is the
weighted (then resampled) particle set, so it can take any shape instead of only a Gaussian.

\vspace{0.5cm}

% ------------------------------------------------------------------------
\subsection*{Q2.}
% ------------------------------------------------------------------------
\noindent\textit{For the measurement step in the particle filter, we used the histogram of the
two patches (the patch around the new particle and the original patch of the object).}

\medskip
\noindent\textit{a. Why do we use the histogram to score the patches? what are the pros and cons
of using this method?}

\medskip
\noindent\textbf{Answer.}

A histogram captures the colour distribution of the patch and deliberately ignores where each pixel sits.

\begin{itemize}[nosep]
    \item \textbf{Pros:} robust to small shifts, rotation, pose/shape change and partial occlusion (a person is non-rigid, so this matters); cheap to compute and compare; gives a compact, normalized appearance signature.
    \item \textbf{Cons:} throws away all spatial / structural information --- two different objects with the same colours look identical; sensitive to illumination and colour changes; background pixels that fall inside the box pollute the histogram; quantizing to 4 bits loses fine colour detail.
\end{itemize}

\vspace{0.5cm}

\noindent\textit{b. We took the patches and computed the histogram and their distances. Why
didn't we use the SSD (sum square difference) between the patches (in the image space)?}

\medskip
\noindent\textbf{Answer.}
% TODO: write your answer here.

  SSD compares pixels position-by-position, so it assumes the two patches are perfectly aligned. A one-pixel shift, a small rotation, a scale change or any deformation makes SSD explode — even for the same object.
\vspace{0.5cm}

\noindent\textit{c. Can you think of another method to compare the patches? What are the pros
and cons of the new method?}

\medskip
\noindent\textit{You can describe an example if it helps you.}

\medskip
\noindent\textbf{Answer.}
% TODO: write your answer here.

  As we seen in Q2.a, one of the cons for using
\vspace{0.5cm}

% ------------------------------------------------------------------------
\subsection*{Q3.}
% ------------------------------------------------------------------------
\noindent\textit{Will particle filter work when the tracked object changes its scale (like a
person walking towards and away from the camera)? What about the viewpoint of the object (like
a person dancing)?}

\medskip
\noindent\textbf{Answer.}
% TODO: write your answer here.

\vspace{0.5cm}

% ------------------------------------------------------------------------
\subsection*{Q4.}
% ------------------------------------------------------------------------
\noindent\textit{Can you think about a way to know when to update the template? Think of a
scenario where we track a person. We keep a patch of the object's appearance in the first frame,
and then compare each new patch to it. One of the problems is that, as time passes, the original
patch no longer represents the object (light, viewpoint, \ldots). Can you think of a way to
determine when to update the patch of the object? And how to determine the new patch (based on
the previous frame/ based on the last two frames/\ldots)?}

\medskip
\noindent\textbf{Answer.}
% TODO: write your answer here.

\vspace{0.5cm}

\end{document}